\scanpage{46}

Considérons la fonction polynomiale sur $\underline P_1$

\[
x\longmapsto (\underbrace{x-x'}_{\substack{\rotatebox{90}{$\scriptstyle\ni$}\\ L}})(\underbrace{x-x''}_{\substack{\rotatebox{90}{$\scriptstyle\ni$}\\ L}})\in L^{\otimes 2}
\]

\[
\text{i.e.}\qquad x\longmapsto x^2-Tx\deleted{\unsure{+N|\psi(x)}}+N
\]

restriction de la forme quadratique

\[
Q:\quad x\longmapsto x^2-\inserted{\psi(x)}Tx\deleted{\unsure{Q(x)|\psi(x)}}+\deleted{\unsure{N}}\psi(x)^2N,\qquad E\longrightarrow L^{\otimes 2}\subset\operatorname{Sym}^2(E)
\]

sur $E$. Cette forme est définie \inserted{par cette formule} en tous les cas, \unsure{supposée|supposons|sans supposer} existence de $x',x''$), elle prend ses valeurs dans $\operatorname{Sym}^2(E)$, mais \deleted{on} vérifie qu'elle est à valeurs dans $L^{\otimes 2}\subset\operatorname{Sym}^2$ (argument ``universel''\ldots) Sa restriction à $L$ est

\[
x\longmapsto Q(x)=x^2:\quad L\longrightarrow L^{\otimes 2}
\]

(partie ``hom.'' de la fonction polynomiale \inserted{de degré 2} envisagée sur $\underline P_1$).

\uline{NB} Equiv. de cat. fibrées

\[
\begin{aligned}
\substack{\text{rev. quadratiques}\\\text{de } S}
&\simeq\left\{\begin{aligned}
&\deleted{\text{Triples}}\ \text{Triples }(\deleted{\unsure{\underline P_1}}\underline L,\underline P_1,Q)\text{ sur }S\\
&\text{a) }\underline L\ \text{Module inv.}\\
&\text{b) }\underline P_1\quad \underline L\text{-torseur}\\
&\text{c) }Q:P_1\longrightarrow L^{\otimes 2}\quad\text{fonct. pol. de degré }2\\
&\text{à }\deleted{\text{partie}}\text{ homogène ``}x\longmapsto x^2\text{''}:\underline L\longrightarrow\underline L^{\otimes 2}
\end{aligned}\right.
\end{aligned}
\]

($Q$ est section d'un torseur sous $\deleted{A}\underline{\operatorname{Hom}}(\underset{\substack{\text{module envelop-}\\\text{pant de }P_1}}{E},\underline L^{\otimes 2})\simeq\mathscr A\otimes\underline L^{\otimes 2}$~; ext. de $\underline L$ par $\underline L^{\otimes 2}$~* \ldots).

% slanted
\margin{* Si $T$ $(=x'+x'')$ est donné d'avance (dans $\Gamma\underline P_2$), alors l'indétermination est \uline{seulement} du $N$, de façon que $x^2+Tx+N$ soit $\in\Gamma\underline L^{\otimes 2}$, \inserted{\unsure{donc indéterminé}} \unsure{donc} mod $\Gamma L^{\otimes 2}$. Si $2$ inv. sur $\mathscr O_S$, alors cette indétermination est levée par la \unsure{considération|donnée} de $\delta\in\Gamma\underline L^{\otimes 2}$, \unsure{vu la} \unsure{form.}~: $4N=T^2-\delta$ Mais il \deleted{\unsure{faut}} faut $T^2\equiv\delta\ (4)$}

Le rev\textsuperscript{t} quadr. associé à ces données peut être interprété comme le ``schéma des zéros de $Q$'', interprété lui-$=$ comme section du faisceau structural du fibré affine $P_1$ correspondant à $\underline P_1$.

% archive sheet number: 45
